\chapter{Hyperhidrosis}

\section*{Definition and Overview}
Hyperhidrosis is a condition of excessive sweating beyond what is needed to regulate body temperature. It most commonly affects the palms, soles, armpits, and face, and it can occur in response to heat or exercise or without any trigger at any time. Primary hyperhidrosis usually begins in childhood or adolescence and is not caused by another condition, while secondary hyperhidrosis is caused by an underlying medical condition or medication. Hyperhidrosis is harmless to physical health but can significantly affect quality of life, and effective treatments are available.

\section*{Epidemiology}
Primary hyperhidrosis affects around 1--3\% of the population, making it more common than is often appreciated. It usually begins in adolescence and often runs in families. The palms and soles are most commonly affected, followed by the armpits. Many people do not seek help, embarrassed by their symptoms, and the condition is associated with anxiety, social withdrawal, and reduced work performance. With treatment, most people gain good control of their sweating.

\section*{Aetiology and Pathophysiology}
\begin{itemize}
    \item \textbf{Primary hyperhidrosis:} overactivity of the sympathetic nervous system causes the sweat glands to work excessively; the exact cause is unknown, but it runs in families
    \item \textbf{Secondary hyperhidrosis:} caused by underlying conditions, including thyroid disease, diabetes, menopause, infections, and some cancers
    \item \textbf{Medications:} some antidepressants and other drugs cause sweating
    \item \textbf{Pathophysiology:} the sweat glands themselves are normal; the problem is excessive nerve stimulation
    \item \textbf{Sites:} eccrine glands in the palms, soles, armpits, and face
    \item \textbf{Triggers:} heat, exercise, stress, and anxiety worsen sweating
\end{itemize}

\section*{Clinical Features}
\begin{itemize}
    \item Visible, excessive sweating of the palms, soles, armpits, or face
    \item Sweating that occurs without heat or exercise
    \item Sweating that interferes with daily activities, work, or social life
    \item Skin problems: maceration, fungal infections, and odour
    \item Difficulty gripping objects or writing
    \item The need to change clothes or shower repeatedly
    \item Night sweats may indicate a secondary cause
\end{itemize}

\section*{Differential Diagnosis}
\begin{itemize}
    \item Primary hyperhidrosis is diagnosed when no underlying cause is found
    \item Secondary causes include thyroid disease, diabetes, menopause, infection, and malignancy
    \item Medication-related sweating
    \item Normal sweating, which is proportionate to heat and activity
    \item Investigations exclude secondary causes, especially in adult-onset or whole-body sweating
\end{itemize}

\section*{When to Seek Medical Care}
See a doctor if sweating is excessive, affects your quality of life, or is accompanied by other symptoms such as weight loss, fever, or night sweats.

\textbf{Call 000 (or local emergency services) for:}
\begin{itemize}
    \item Drenching sweats with chest pain, breathlessness, or collapse
    \item Fever with confusion or severe illness
    \item Signs of hypoglycaemia or other acute emergencies
    \item Any life-threatening emergency
\end{itemize}

\section*{Diagnosis and Investigations}
\begin{itemize}
    \item \textbf{Clinical assessment:} the pattern, sites, and triggers of sweating
    \item \textbf{History:} onset, family history, medications, and associated symptoms
    \item \textbf{Blood tests:} thyroid function, blood sugar, and other tests as indicated
    \item \textbf{Assessment of impact:} the effect on work, social life, and wellbeing
    \item \textbf{Referral:} to dermatology or specialist sweating clinics for severe cases
\end{itemize}

\section*{Management}
\begin{itemize}
    \item \textbf{Antiperspirants:} aluminium chloride antiperspirants are the first-line treatment
    \item \textbf{Iontophoresis:} passing a mild current through water applied to hands and feet
    \item \textbf{Botox injections:} injections into the armpits or other sites block the nerves that stimulate sweat glands
    \item \textbf{Medication:} anticholinergic tablets reduce sweating in some people
    \item \textbf{Microwave and laser treatments:} destroy sweat glands in the armpits
    \item \textbf{Surgery:} endoscopic thoracic sympathectomy for severe hand sweating, with careful consideration of risks
    \item \textbf{Treatment of underlying causes:} for secondary hyperhidrosis
\end{itemize}

\section*{Complications}
\begin{itemize}
    \item Skin infections, including fungal infections and maceration
    \item Body odour and social embarrassment
    \item Anxiety, depression, and social withdrawal
    \item Impaired work and education performance
    \item Treatment side effects, including compensatory sweating after surgery
    \item Dehydration in hot conditions with excessive whole-body sweating
\end{itemize}

\section*{Prognosis and Outcomes}
Hyperhidrosis is a chronic condition, but it responds well to treatment. Most people achieve good control with antiperspirants or iontophoresis, and Botox is highly effective for armpit sweating, with effects lasting months. Surgery helps severe hand sweating in most cases, though compensatory sweating elsewhere can occur. With the range of treatments available, most people with hyperhidrosis can manage their symptoms and live normally.

\section*{Prevention and Self-Care}
\begin{itemize}
    \item Use antiperspirant rather than deodorant, applied to dry skin at night
    \item Wear breathable, loose clothing
    \item Use sweat-absorbing insoles and moisture-wicking socks
    \item Manage stress and anxiety, which trigger sweating
    \item Stay cool and hydrated
    \item Treat skin problems early to prevent infection
    \item Seek treatment early rather than suffering in silence
    \item Discuss all treatment options with your doctor
\end{itemize}

\section*{Sources and Further Reading}
The following reputable sources provide further information for healthcare students and patients seeking reliable, current advice about hyperhidrosis.

\begin{itemize}
    \item Healthdirect. \textit{Hyperhidrosis (excessive sweating).} https://www.healthdirect.gov.au/hyperhidrosis
    \item Australasian College of Dermatologists. \textit{Hyperhidrosis information.} https://www.dermcoll.edu.au/
    \item International Hyperhidrosis Society. \textit{Information and treatment options.} https://www.sweathelp.org/
    \item Better Health Channel. \textit{Excessive sweating.} https://www.betterhealth.vic.gov.au/
    \item NHS. \textit{Hyperhidrosis.} https://www.nhs.uk/conditions/hyperhidrosis/
    \item Mayo Clinic. \textit{Hyperhidrosis: symptoms and causes.} https://www.mayoclinic.org/
\end{itemize}
